% ------------------------------------------------------------------------------
% Chapter 6
% ------------------------------------------------------------------------------
\chapter{Conclusion} % enter the name of the chapter here
\label{cha:ch6_conclusion} % enter the chapter label here (for cross referencing)

\section{Introduction}
\label{sec:Ch6 Introduction}
This final chapter concludes the dissertation. It summarises the study, states the conclusions reached in relation to the research hypotheses, sets out the study's contribution to the body of knowledge, offers recommendations for IT leadership practice, restates the study's limitations, proposes directions for future research, and closes with a set of concluding remarks.

\section{Summary of the Study}
\label{sec:Summary of the Study}
This study examined the role of organisational culture and leadership in realising the return on investment of AI coding assistants in software engineering teams, and, in doing so, the role of organisational readiness, team size, developer experience, and adoption duration as contextual factors. A quantitative, post-positivist, cross-sectional design was adopted (Chapter Three), combining preliminary Pearson correlation screening with a partial least squares structural equation model (PLS-SEM), estimated using a Python-based implementation of the PLS algorithm. Data were collected from 354 software engineering practitioners, predominantly software developers, senior developers, and technical leads, in large South African organisations using AI coding assistants. The structural model tested four hypotheses concerning the relationships between organisational culture, AI leadership and governance support, the assimilation quality of AI coding assistants, and ROI, net of organisational readiness, team size, developer experience, and adoption duration (Chapter Four), and the findings were interpreted in relation to the literature in Chapter Five.

\section{Conclusions on the Research Hypotheses}
\label{sec:Conclusions on the Research Hypotheses}
Table~\ref{tab:conclusions_summary} maps each hypothesis to its outcome.

\begin{table}[htbp]
	\centering
	\caption{Summary of conclusions by hypothesis}
	\label{tab:conclusions_summary}
	\resizebox{\linewidth}{!}{%
		\begin{tabular}{lp{7.7cm}l}
			\hline
			\textbf{Hyp.} & \textbf{Statement} & \textbf{Outcome} \\
			\hline
			H1 & Organisational culture $\leftrightarrow$ ROI (total effect) & Supported (indirect only) \\
			H2 & AI leadership and governance support $\leftrightarrow$ ROI (total effect) & Not supported \\
			H3 & Assimilation quality $\leftrightarrow$ ROI (direct effect, $\beta=.356$) & Supported \\
			H4a & Assimilation quality mediates culture $\rightarrow$ ROI & Supported (indirect-only mediation) \\
			H4b & Assimilation quality mediates leadership $\rightarrow$ ROI & Not supported \\
			\hline
		\end{tabular}%
	}
\end{table}

On the basis of these results, it is concluded that organisational culture is associated with the ROI realised from AI coding assistants, but that this association operates entirely through the quality with which the tools are assimilated into software engineering practice, rather than through any additional direct channel. AI leadership and governance support, by contrast, did not show an independent, statistically supported association with ROI once organisational culture was accounted for in the same model, a result attributable to the substantial empirical overlap between the leadership and culture measures rather than to leadership being unimportant in an absolute sense. Assimilation quality is confirmed as a genuine, direct contributor to ROI, and organisational readiness, entered as a control variable rather than a hypothesised predictor, emerged as the single strongest predictor of ROI in the entire model, ahead of assimilation quality itself. Read together, these results reposition assimilation quality as an important but not solitary driver of ROI, and elevate organisational readiness, technical infrastructure, developer skills, governance mechanisms, and value-assessment capability, to a position of practical and theoretical prominence that the original hypothesised model did not anticipate.

\section{Answering the Research Questions and Objectives}
\label{sec:Answering the Research Questions}
The study was structured around one main research question and five sub-research questions (Section~\ref{sec:Research questions}), corresponding to the five research objectives set out in Section~\ref{sec:Research aim and objectives}. This section states each in turn and summarises how the study addressed it.

\textbf{Main research question: How do organisational culture and leadership shape the ROI of AI coding assistants in South African software engineering teams?} Organisational culture shapes ROI, but does so entirely indirectly, through the quality with which AI coding assistants are assimilated into engineering practice (H1, H4a), rather than through any direct channel. AI leadership and governance support did not demonstrate an independent shaping effect on ROI once organisational culture was accounted for in the same model (H2, H4b), a result attributable to substantial empirical overlap between the two constructs rather than to leadership being irrelevant. Organisational readiness, entered as a control rather than as part of the original hypothesised answer to this question, in fact shaped ROI more strongly than either culture or leadership did.

\textbf{RQ1: What financial costs and benefits define ROI for AI coding assistants in software engineering teams?} This question was addressed through measurement rather than hypothesis testing. ROI was operationalised in Chapter Three as perceived business value across three levels, task-level productivity and quality effects, workflow-level integration and delivery effects, and firm-level cost, speed-to-market, innovation, and financial-value effects, using a 15-item scale. The reliability and measurement-model results in Chapter Four (Cronbach's alpha of .990 for the overall scale, outer loadings of .909 to .962, AVE of .880) confirm that these 15 items cohere as a single, well-measured perceived-ROI construct, while the discriminant validity results (Section~\ref{sec:Discriminant validity}) show that the three levels could not be empirically separated from one another in this sample, so RQ1 is answered at the level of an overall perceived-ROI construct rather than three fully distinct financial and non-financial sub-constructs.

\textbf{RQ2: How do organisational culture attributes influence productivity and quality outcomes from AI coding assistants?} Organisational culture was significantly associated with ROI (which incorporates the task- and workflow-level productivity and quality items forming part of RQ1's operationalisation) at the preliminary stage, and the structural model shows this influence operates specifically by improving assimilation quality (H1, H4a): culture's direct effect on ROI was not significant once assimilation quality was included, but its indirect effect through assimilation quality was ($\beta = .206$, 95\% CI [.025, .409]).

\textbf{RQ3: Which leadership behaviours strongly affect ROI outcomes from AI coding assistants?} As discussed in Chapter Five, this study's evidence for RQ3 is qualified. The bundled leadership construct did not show an independent structural effect on ROI (H2). The supplementary item-level analysis in Section~\ref{sec:Leadership Behaviours Supplementary Analysis} found governance clarity ($r = .431$) and executive sponsorship ($r = .410$) to have the strongest individual bivariate associations with ROI, ahead of resource support ($r = .360$) and metric alignment ($r = .344$), but this ordering did not survive, and could not be formally tested, once culture was controlled for. RQ3's answer is therefore that no leadership behaviour, individually or bundled, demonstrated an effect on ROI independent of organisational culture in this sample.

\textbf{RQ4: How can financial and organisational factors be integrated into a practical ROI framework for large South African organisations?} The validated structural model itself is this study's answer to RQ4: a PLS-SEM framework integrating organisational culture, AI leadership and governance support, assimilation quality, organisational readiness, and three demographic controls, which explained 81.6\% of the variance in perceived ROI ($R^2 = .816$) and showed acceptable approximate fit (SRMR $= .054$). The framework's practical contribution is the sequencing it implies: organisational readiness and assimilation quality are the two constructs with a confirmed, independent, and directly actionable relationship with ROI, and should accordingly anchor a practical ROI framework, with culture and leadership treated as upstream enablers of assimilation quality rather than as independent levers. This sequencing follows directly from the importance-performance map analysis reported in Chapter Four (Section~\ref{sec:IPMA}), which places organisational readiness uniquely in the map's priority-for-action quadrant, combining the highest importance to ROI with the lowest relative performance of the four constructs examined.

\textbf{RQ5: What practical guidance can be offered to managers to maximise returns from AI coding assistant investments through cultural and leadership interventions?} This is addressed in full in Section~\ref{sec:Recommendations for IT Leadership Practice}; in summary, managers are advised to prioritise organisational readiness first, assimilation quality second, and to treat culture and leadership interventions as complementary, mutually reinforcing investments rather than as separable levers, given that this study could not distinguish their independent contributions.

\section{Contribution to Knowledge}
\label{sec:Contribution to Knowledge}
\subsection{Theoretical Contribution}
\label{sec:Theoretical Contribution}
This study adds to the literature on generative AI adoption and IT leadership in three ways. First, it empirically tests, rather than merely proposes, a mediating role for assimilation quality between organisational culture and ROI, and shows that this mediation is best classified as indirect-only: culture's contribution to ROI is fully channelled through assimilation quality. Second, by modelling AI leadership and governance support alongside organisational culture in a single structural model, rather than treating each as an isolated predictor, the study demonstrates empirically how two theoretically distinct but empirically overlapping constructs can produce misleading conclusions when tested only through separate bivariate associations, a methodological caution of relevance beyond this study's specific constructs. Third, by including organisational readiness, team size, developer experience, and adoption duration as controls rather than omitting them, the study surfaces organisational readiness as an unexpectedly dominant predictor of ROI, suggesting that future theoretical models of AI-driven ROI in software engineering should consider foundational organisational readiness as a first-order construct rather than a background control variable.

\subsection{Practical Contribution}
\label{sec:Practical Contribution}
For an IT leader in a large organisation, this study provides evidence-based justification for a specific sequencing of effort: organisational readiness, ensuring the technical infrastructure, developer skills, governance mechanisms, and value-assessment capability are in place, was the strongest single predictor of perceived ROI from AI coding assistants in this sample, ahead of the more behavioural, day-to-day quality of assimilation on which much practitioner attention tends to focus. Investment in culture remains worthwhile, but this study indicates that its value is realised specifically by improving assimilation quality, rather than through some other, more direct route; and clarity and support activity, that is separable and independently effective from a supportive culture.

\section{Recommendations for IT Leadership Practice}
\label{sec:Recommendations for IT Leadership Practice}
Reflecting the dissertation's focus on the role of IT leadership in realising ROI from AI coding assistants, and reflecting the relative empirical support behind each recommendation, the following are proposed:
\begin{enumerate}
	\item IT leaders should prioritise organisational readiness, technical infrastructure, developer skills, governance mechanisms, and the capability to assess whether AI tools are generating measurable value, as a foundational investment area, given that it was the strongest and most robust predictor of perceived ROI in this study, and that the importance-performance map analysis (Chapter Four, Section~\ref{sec:IPMA}) formally identifies it as the single highest-priority construct for management action.
	\item IT leaders should treat assimilation quality, not tool provisioning or usage volume alone, as a second key indicator to monitor and manage, given its statistically supported direct association with ROI (H3) and its role as the specific pathway through which organisational culture's influence operates (H4a).
	\item Investment in a supportive organisational culture, experimentation, learning orientation, knowledge sharing, and digital openness, should continue, on the understanding that, on this study's evidence, its contribution to ROI operates through assimilation quality rather than through a separate, direct channel.
	\item Executive sponsorship, metric alignment, resource support, and governance clarity should not be deprioritised on the strength of this study's findings alone; rather, organisations should be aware that this study could not separate leadership's independent contribution to ROI from that of organisational culture, and should treat the two as working together rather than assume either operates alone.
	\item Organisations should complement self-reported perceptions of ROI, culture, leadership, and readiness with more objective delivery, quality, and cost metrics where feasible, given the very high, uniformly positive self-report scores, the common method bias identified in this study, and the demonstrated sensitivity of the key estimates to this response pattern.
\end{enumerate}

\section{Limitations of the Study}
\label{sec:Ch6 Limitations}
The limitations set out in Section~\ref{sec:Research Limitations} should be borne in mind when interpreting these conclusions. In particular, the cross-sectional design cannot establish causality with the same confidence as an experimental or longitudinal design; all constructs were measured by self-report from the same respondents, and common method bias was assessed with two converging procedures, Harman's single-factor test (a single unrotated factor accounted for 65.4\% of total variance) and a full-collinearity VIF test (five of eight constructs exceeded the conservative 3.3 threshold), both indicating that common method bias is a genuine concern; a dedicated sensitivity analysis showed that 69.8\% of respondents selected the most positive response option on every substantive item, a response pattern whose exclusion materially weakens, though does not reverse, the estimated strength of the assimilation-quality pathway; organisational culture and AI leadership and governance support were found to be empirically difficult to distinguish from one another on the more stringent HTMT criterion (HTMT = .881, 95\% CI upper bound = .967), although not on the more conservative Fornell-Larcker criterion, which limits the confidence with which their independent effects on ROI can be separated; the three ROI sub-scales were found not to be empirically distinguishable and were therefore represented as a single overall construct, meaning the study cannot speak to whether the predictors relate differently to ROI at the task, workflow, or firm level; the achieved sample, while exceeding both the ten-times-rule and inverse-square-root minimum sample sizes for the paths found to be statistically significant (Section~\ref{sec:Population and Sampling}), was drawn using purposive and stratified sampling from large South African organisations and included a broader mix of engineering roles than originally targeted; and ROI was operationalised through perceptual rather than direct financial indicators. Two further limitations concern the scope of the structural analysis itself, rather than the data: this study did not test for measurement invariance or conduct a multigroup comparison across role or team size (MICOM and PLS-MGA), which the achieved sample would support and which is the most natural extension of the present model; and it did not model unobserved heterogeneity, nonlinear effects, or endogeneity, each a legitimate, more advanced extension beyond the scope of a single master's-level study rather than an oversight.

\section{Recommendations for Future Research}
\label{sec:Recommendations for Future Research}
\begin{enumerate}
	\item A longitudinal or quasi-experimental design would strengthen causal inference beyond what the cross-sectional design used here can support, particularly for the mediating role of assimilation quality and for the unexpectedly dominant role of organisational readiness.
	\item Future studies should incorporate objective delivery, quality, or financial metrics alongside self-reported measures, and should consider collecting predictor and outcome data from different sources or at different time points, to reduce the common method variance and ceiling effects identified in this study.
	\item Because organisational culture and AI leadership and governance support could not be clearly distinguished in this sample, future instrument development should revisit the conceptual and item-level separation of these two constructs, or explicitly model their overlap, before drawing conclusions about their independent effects.
	\item Given the prominence of organisational readiness in this study's results, future research should treat it as a first-order theoretical construct, potentially examining whether it acts as an antecedent, moderator, or precondition for assimilation quality, rather than including it only as a control variable.
	\item Given the very high intercorrelation between the task-, workflow-, and firm-level ROI sub-scales, future instrument development should examine whether these are better treated as a single ROI construct or refined to capture more clearly differentiated levels of value.
	\item Future research should extend the sample beyond large South African organisations, and beyond the broader mix of software engineering roles captured here, to test the generalisability of these findings.
	\item Future research should test for measurement invariance and conduct a multigroup comparison (MICOM and PLS-MGA) across role and team size, which the achieved sample in this study would have supported but which was outside this study's scope, to establish whether the relationships identified here hold equally across different kinds of software engineering teams.
	\item Future research should extend the structural model to account for unobserved heterogeneity, nonlinear effects, and endogeneity, each a legitimate methodological extension that a single master's-level study was not scoped to include.
\end{enumerate}

\section{Concluding Remarks}
\label{sec:Concluding Remarks}
This study set out to examine the role of organisational culture and leadership in realising ROI from AI coding assistants, and finds a more nuanced picture than initially hypothesised. Organisational culture matters, but principally through a single, actionable channel: how well the tools are assimilated into the everyday practice of software engineering teams. AI leadership and governance support, measured separately, did not demonstrate an independent effect once its substantial overlap with culture was accounted for. Most notably, the foundational readiness of the organisation, its infrastructure, skills, governance, and capability to assess value, proved to be the strongest predictor of perceived ROI of all, ahead of the behavioural, assimilation-focused factors that received the greater share of this study's theoretical attention. For IT leaders, the implication is that building a supportive culture and providing well-governed leadership sponsorship remain necessary, but that neither substitutes for ensuring the organisation is genuinely ready, technically, and in terms of skills and governance, to support AI coding assistants in the first place; the return on that investment is realised, or lost, at the intersection of readiness and assimilation, not at either alone.